\chapter{Geltungsbereich und Charakterisierungen}
\label{b12guide:ch:foundations}

Dieser Leitfaden ergänzt Band~12 um eine ausgewählte Übersicht
zentraler Definitionen und Sätze über \emph{Halbordnungen}.
Vorausgesetzt werden ausschließlich Reflexivität, Antisymmetrie
und Transitivität auf einer beliebigen Menge. Der Träger darf leer
sein. Mengentheoretische Grundlage ist ZF ohne zusätzliches
Auswahlaxiom. Der Leitfaden enthält keine Beweise oder Beweisskizzen.

\section{Gegenstand und Bandgrenzen}

Aufgenommen sind Charakterisierungen, die strenge Schreibweise,
Dualität, Extremalbegriffe, ordnungsbezogene Teilmengen sowie
Konstruktionen und das Erzeugen von Halbordnungen. Bedingungen
an einzelne Elemente oder Teilmengen werden jeweils ausdrücklich
genannt; sie werden nicht als zusätzliche Axiome an alle
Halbordnungen vorausgesetzt.

Allgemeine Sätze über Funktionen sowie Injektivität, Surjektivität
und Bijektivität bleiben den Bänden~5--8 vorbehalten.
Äquivalenzrelationen und Klassenbildung gehören zu Band~11;
auch die Präordnungen und ihre Quotienten am Ende des Hauptbandes~12
werden hier nicht behandelt.
Schranken, Infima und Suprema, Paarinfima und Paarsuprema,
totale Ordnungen, Wohlordnungen und Auswahlaxiome sowie
Halbverbände und Verbände bleiben den Bänden~13--16 und~45
vorbehalten. Sätze, die Endlichkeit voraussetzen, sind ebenfalls
nicht Gegenstand dieses Leitfadens.

Die strenge Ordnung ist eine gleichwertige Schreibweise derselben
Halbordnung. Unter- und Obermengen sowie Antiketten werden als
Teilmengen einer Halbordnung betrachtet. Eine Theorie von
Ordnungsabbildungen wird nicht entwickelt.

Die Zusammenstellung ist eine Referenzliste wichtiger bekannter
Resultate, keine vollständige Liste aller bekannten oder ableitbaren
Aussagen. Die Nummern mit Präfix $12H$ gehören ausschließlich
zu dieser Ergänzung.

\section{Grundbegriff und Relationsnotation}

\begin{definition}[Halbordnung]\label{b12guide:def:order}
Eine \emph{Halbordnung} oder \emph{partielle Ordnung} auf einer
Menge $P$ ist ein binärer Relationsoperator $\leq$, für den
für alle $x,y,z\in P$ gilt:
\[
\begin{aligned}
 &x\leq x &&\text{(Reflexivität)},\\
 &(x\leq y\ \land\ y\leq x)\ \Rightarrow\ x=y
     &&\text{(Antisymmetrie)},\\
 &(x\leq y\ \land\ y\leq z)\ \Rightarrow\ x\leq z
     &&\text{(Transitivität)}.
\end{aligned}
\]
Das Paar $(P,\leq)$ heißt \emph{halbgeordnete Menge} oder
\emph{Poset}. Im folgenden ist $(P,\leq)$ stets eine
Halbordnung, sofern nicht ausdrücklich eine beliebige Relation
untersucht wird.
\end{definition}

Für Aussagen über Relationsmengen verwenden wir den
\emph{Relationsgraphen}
\[
 R=\{(x,y)\in P\times P\mid x\leq y\},\qquad
 \Delta_P=\{(x,x)\mid x\in P\}.
\]
Die Menge aller Graphen von Halbordnungen auf $P$ heiße
$\PO(P)$. Zu einer Paarmenge $R\subseteq P\times P$
bezeichnet $\leq_R$ den Operator mit
$x\leq_R y\Longleftrightarrow(x,y)\in R$.
Die Schreibweise $R\subseteq S$ bedeutet, daß jede in $R$
geltende Vergleichsbeziehung auch in $S$ gilt.
Weiter sei
\[
\begin{aligned}
 R^{\mathrm{op}}&=\{(y,x)\mid(x,y)\in R\},\\
 R\circ R&=\{(x,z)\in P\times P\mid
       \exists y\in P\;((x,y)\in R\land(y,z)\in R)\}.
\end{aligned}
\]

\begin{theorem}[Charakterisierung durch Relationsgleichungen]
\label{b12guide:thm:relation-criterion}
Für jede Relation $R\subseteq P\times P$ gilt
\[
 R\in\PO(P)\quad\Longleftrightarrow\quad
 R\cap R^{\mathrm{op}}=\Delta_P\ \land\ R\circ R=R.
\]
\end{theorem}

\begin{theorem}[Lokales Kriterium auf höchstens drei Elementen]
\label{b12guide:thm:local}
Eine Relation $R\subseteq P\times P$ ist genau dann eine
Halbordnung, wenn für jede Teilmenge $B\subseteq P$ mit
höchstens drei Elementen die Relation
$R\cap(B\times B)$ eine Halbordnung auf $B$ ist.
Eine Endlichkeit von $P$ wird nicht vorausgesetzt.
\end{theorem}

\begin{theorem}[Diskrete Ordnung und Randfälle]
\label{b12guide:thm:discrete}
Die Gleichheitsrelation $\Delta_P$ ist eine Halbordnung auf jeder
Menge $P$, und $\Delta_P\subseteq R$ für jedes $R\in\PO(P)$.
Sie heißt \emph{diskrete Ordnung}.
Auf der leeren Menge und auf einer einelementigen Menge ist
sie die einzige Halbordnung. Ferner gilt
\[
 P\times P\in\PO(P)\quad\Longleftrightarrow\quad
 P\text{ hat höchstens ein Element}.
\]
\end{theorem}

